% kap04.tex — Poisson-parenteser og kanonisk struktur
% Hamilton-Jacobi-kompendiet

\chapter{Poisson-parenteser og kanonisk struktur}
\label{kap:poisson}

\section{Tidsudviklingen af en vilkårlig størrelse}

I kapitel~\ref{kap:hamilton} udledte vi, at $dH/dt=\partial H/\partial t$
langs en faktisk bane. Det naturlige næste spørgsmål er: hvad gælder
for tidsafledningen af en \emph{vilkårlig} funktion
$f(q,p,t)$ af faserumsvariablene — ikke kun $H$ selv? Svaret på dette
spørgsmål fører os til at indføre \emph{Poisson-parentesen}\index{Poisson-parentes}, som
bliver det centrale algebraiske redskab for resten af kompendiet, og
som — det skal siges allerede her — er den klassiske forløber for
kommutatoren i kvantemekanikken.

Ved kædereglen og Hamiltons ligninger \eqref{eq:hamiltons-ligninger}:
\begin{equation}
  \dv{f}{t} = \sum_i\left(\pdv{f}{q_i}\dot q_i + \pdv{f}{p_i}\dot p_i\right) + \pdv{f}{t}
  = \sum_i\left(\pdv{f}{q_i}\pdv{H}{p_i} - \pdv{f}{p_i}\pdv{H}{q_i}\right) + \pdv{f}{t}.
  \label{eq:df-dt-udledning}
\end{equation}
Summen på højresiden — uden $\partial f/\partial t$-leddet — optræder
så ofte og med en så bestemt algebraisk struktur, at den fortjener sit
eget symbol og navn.

\section{Definition af Poisson-parentesen}

For to funktioner $f(q,p,t)$ og $g(q,p,t)$ på faserummet defineres
\emph{Poisson-parentesen}
\begin{equation}
  \boxed{\{f,g\} \equiv \sum_{i=1}^n\left(\pdv{f}{q_i}\pdv{g}{p_i} - \pdv{f}{p_i}\pdv{g}{q_i}\right).}
  \label{eq:def-poisson}
\end{equation}
Med denne definition bliver \eqref{eq:df-dt-udledning} til den
kompakte og centrale formel
\begin{equation}
  \boxed{\dv{f}{t} = \{f,H\} + \pdv{f}{t}.}
  \label{eq:bevaegelse-poisson}
\end{equation}
Dette er \emph{bevægelsesligningen for en vilkårlig faserumsfunktion},
udtrykt via Poisson-parentesen med $H$. Vi kan endda genskrive
Hamiltons ligninger selv i denne form: sæt $f=q_i$ og $f=p_i$
(begge er trivielt tidsuafhængige som funktioner, $\partial q_i/\partial t=0$),
og brug definitionen \eqref{eq:def-poisson} direkte:
\begin{equation}
  \dot q_i = \{q_i,H\} = \pdv{H}{p_i}, \qquad
  \dot p_i = \{p_i,H\} = -\pdv{H}{q_i},
\end{equation}
i fuld overensstemmelse med \eqref{eq:hamiltons-ligninger}. Hamiltons
ligninger er altså blot specialtilfælde af den generelle formel
\eqref{eq:bevaegelse-poisson}.

\section{De fundamentale parenteser}

Ved direkte indsættelse i \eqref{eq:def-poisson}, med $q_i$ og $p_i$
som de indsatte funktioner, og med Kronecker-deltaet
$\delta_{ij}=1$ for $i=j$ og $0$ ellers (fordi $\partial q_i/\partial q_j=\delta_{ij}$
og $\partial p_i/\partial p_j=\delta_{ij}$, mens $q_i$ og $p_j$ er
uafhængige, så $\partial q_i/\partial p_j=0$), finder man
\begin{equation}
  \boxed{\{q_i,q_j\} = 0, \qquad \{p_i,p_j\} = 0, \qquad \{q_i,p_j\} = \delta_{ij}.}
  \label{eq:fundamentale-parenteser}
\end{equation}
Disse tre relationer kaldes de \emph{fundamentale
Poisson-parenteser}\index{Poisson-parentes!fundamentale}.
De er selve definitionen af, hvad det betyder for et sæt variable
$(q_i,p_i)$ at være \emph{kanonisk konjugerede}\index{kanonisk konjugerede variable} — en egenskab, vi
udnytter systematisk, når vi i kapitel~\ref{kap:kanoniske-transf}
spørger, hvilke koordinatskift på faserummet bevarer Hamiltons
ligningers form.

\begin{bemaerkning}
Relationen $\{q_i,p_j\}=\delta_{ij}$ er den klassiske forgænger til
kommutationsrelationen $[\hat q_i,\hat p_j]=i\hbar\,\delta_{ij}$ i
kvantemekanikken (se [[kvantemekanik-kompendiet]]). Reglen for at gå
fra klassisk til kvantemekanik — \emph{kanonisk kvantisering} — er
netop at erstatte Poisson-parenteser med kommutatorer efter reglen
$\{f,g\}\to \frac{1}{i\hbar}[\hat f,\hat g]$. Vi vil ikke udføre
denne kvantisering i dette kompendium (det er det senere QFT-
kompendiums opgave at gøre det systematisk), men det er værd at bære
med sig, at hele den algebraiske struktur, vi bygger op i dette
kapitel, overlever næsten ord for ord ind i kvanteteorien.
\end{bemaerkning}

\section{Algebraiske egenskaber}

Poisson-parentesen har en række egenskaber, som gør den til en såkaldt
\emph{Lie-parentes}\index{Lie-algebra} — den samme abstrakte struktur, man møder i teorien
for kontinuerte symmetrigrupper. Vi angiver dem her; de følger alle
direkte af definitionen \eqref{eq:def-poisson} ved indsættelse (se
opgave 1).

\begin{itemize}
  \item \textbf{Antisymmetri:} $\{f,g\} = -\{g,f\}$. Specielt
        $\{f,f\}=0$.
  \item \textbf{Bilinearitet:} $\{af+bg,h\} = a\{f,h\}+b\{g,h\}$ for
        konstanter $a,b$ (og tilsvarende i det andet argument, via
        antisymmetrien).
  \item \textbf{Leibniz-regel (produktregel):}
        $\{fg,h\} = f\{g,h\} + \{f,h\}g$.
  \item \textbf{Jacobi-identiteten:}\index{Jacobi-identiteten}
        \begin{equation}
          \{f,\{g,h\}\} + \{g,\{h,f\}\} + \{h,\{f,g\}\} = 0.
        \end{equation}
\end{itemize}

Jacobi-identiteten er den mindst intuitive, men også den vigtigste af
de fire: det er den, der garanterer, at Poisson-parenteser (og senere
kommutatorer) definerer en konsistent algebraisk struktur, som man kan
bygge en gruppeteori af symmetritransformationer ovenpå.

\begin{bemaerkning}
Et sæt størrelser, der er lukket under Poisson-parentesen — dvs.\
hvor $\{f,g\}$ altid selv kan skrives som en (lineær) kombination af
størrelserne i sættet — udgør, hvad man kalder en \emph{Lie-algebra}.
Dette er præcis den struktur, der ligger bag begrebet "symmetrigruppe"
i fysikken: banemomentkomponenterne $(L_x,L_y,L_z)$, som vi ser i
næste afsnit, er det arketypiske eksempel, og deres Poisson-
algebra er den klassiske skygge af rotationsgruppens
kommutationsrelationer i kvantemekanikken.
\end{bemaerkning}

\section{Bevarelseslove via Poisson-parenteser}

Formel \eqref{eq:bevaegelse-poisson} giver et umiddelbart og meget
anvendt kriterium for bevarelse: hvis $f$ ikke afhænger eksplicit af
tiden ($\partial f/\partial t=0$), så er
\begin{equation}
  \boxed{f \text{ er bevaret} \quad\Longleftrightarrow\quad \{f,H\}=0.}
  \label{eq:bevarelse-poisson}
\end{equation}
Dette generaliserer både resultatet for cykliske koordinater
(kapitel~\ref{kap:lagrange}) og for $H$ selv
(kapitel~\ref{kap:hamilton}, hvor $\{H,H\}=0$ trivielt pga.\
antisymmetrien, så $H$ er bevaret netop når $\partial H/\partial t=0$,
som vi allerede vidste).

\begin{bemaerkning}
Et vigtigt strukturelt resultat, som følger direkte af Jacobi-
identiteten (se opgave 2), er, at hvis $f$ og $g$ begge er bevarede
størrelser (begge Poisson-kommuterer med $H$), så er deres Poisson-
parentes $\{f,g\}$ \emph{også} en bevaret størrelse. Bevarede
størrelser er altså lukkede under Poisson-parentesen — de udgør netop
en Lie-algebra i den forstand, der blev nævnt ovenfor. Dette er
oprindelsen til, at bevægelseskonstanter for et system med
tilstrækkeligt mange symmetrier ofte "kommer i familier", der danner
en gruppe (rotationsgruppen for banemoment er det stående eksempel).
\end{bemaerkning}

\subsection{Eksempel: banemomentets Poisson-algebra}

\begin{eksempel}
Betragt banemomentets\index{banemoment} tre kartesiske komponenter for en partikel i tre
dimensioner,
\begin{equation}
  L_x = yp_z - zp_y, \qquad L_y = zp_x - xp_z, \qquad L_z = xp_y - yp_x,
\end{equation}
udtrykt ved de kanoniske variable $(x,y,z,p_x,p_y,p_z)$, som opfylder
de fundamentale parenteser \eqref{eq:fundamentale-parenteser} (med
$q_1,q_2,q_3=x,y,z$). Vi beregner $\{L_x,L_y\}$ direkte fra
definitionen \eqref{eq:def-poisson},
\begin{equation}
  \{L_x,L_y\} = \sum_{i\in\{x,y,z\}}\left(\pdv{L_x}{i}\pdv{L_y}{p_i} - \pdv{L_x}{p_i}\pdv{L_y}{i}\right),
\end{equation}
ved først at opskrive alle de nødvendige partielle afledte:
\begin{equation}
  \pdv{L_x}{x}=0,\quad \pdv{L_x}{y}=p_z,\quad \pdv{L_x}{z}=-p_y,\quad
  \pdv{L_x}{p_x}=0,\quad \pdv{L_x}{p_y}=-z,\quad \pdv{L_x}{p_z}=y,
\end{equation}
\begin{equation}
  \pdv{L_y}{x}=-p_z,\quad \pdv{L_y}{y}=0,\quad \pdv{L_y}{z}=p_x,\quad
  \pdv{L_y}{p_x}=z,\quad \pdv{L_y}{p_y}=0,\quad \pdv{L_y}{p_z}=-x.
\end{equation}
Bidraget fra $i=x$ er $\left(\partial_xL_x\right)\left(\partial_{p_x}L_y\right)-\left(\partial_{p_x}L_x\right)\left(\partial_xL_y\right)=0\cdot z-0\cdot(-p_z)=0$,
og bidraget fra $i=y$ er ligeledes $p_z\cdot 0-(-z)\cdot 0=0$: begge
forsvinder, fordi enten $L_x$ eller $L_y$ mangler den relevante
afhængighed. Kun $i=z$ bidrager:
\begin{equation}
  \left(\pdv{L_x}{z}\right)\left(\pdv{L_y}{p_z}\right) - \left(\pdv{L_x}{p_z}\right)\left(\pdv{L_y}{z}\right)
  = (-p_y)(-x) - (y)(p_x) = xp_y - yp_x = L_z.
\end{equation}
Vi finder altså
\begin{equation}
  \boxed{\{L_x,L_y\} = L_z,}
  \label{eq:Lx-Ly}
\end{equation}
og ved cyklisk ombytning af $(x,y,z)$:
\begin{equation}
  \{L_y,L_z\} = L_x, \qquad \{L_z,L_x\} = L_y.
\end{equation}
Banemomentets tre komponenter er altså \emph{lukkede} under Poisson-
parentesen — de danner en Lie-algebra, isomorf til rotationsgruppen
$SO(3)$'s Lie-algebra. Dette er den klassiske skygge af de velkendte
kommutationsrelationer $[\hat L_x,\hat L_y]=i\hbar\hat L_z$ (cyklisk) i
kvantemekanikken (se [[kvantemekanik-kompendiet]]).
\end{eksempel}

\section{Liouvilles sætning: bevarelse af faserumsvolumen}

Vi slutter kapitlet med et geometrisk resultat, der binder Poisson-
strukturen sammen med faserumsbilledet fra
kapitel~\ref{kap:hamilton}. Betragt en samling af systemer med
starttilstande fordelt over et lille område $\Omega_0$ i faserummet
(fx alle med næsten samme $(q,p)$, men med små variationer). Under
tidsudviklingen bevæger hvert punkt i $\Omega_0$ sig efter Hamiltons
ligninger, og området deformeres til et nyt område $\Omega_t$.

\emph{Liouvilles sætning}\index{Liouvilles sætning} siger, at \emph{volumenet} af dette område i
faserummet er \emph{bevaret}:
\begin{equation}
  \text{Vol}(\Omega_t) = \text{Vol}(\Omega_0) \qquad \text{for alle } t.
  \label{eq:liouville}
\end{equation}
Den underliggende grund er, at Hamiltons ligninger definerer et
"strømningsfelt" på faserummet,
$\vb{v}=(\dot q_1,\dots,\dot q_n,\dot p_1,\dots,\dot p_n)$, hvis
divergens er identisk nul:
\begin{equation}
  \sum_i\left(\pdv{\dot q_i}{q_i} + \pdv{\dot p_i}{p_i}\right)
  = \sum_i\left(\pdv{}{q_i}\pdv{H}{p_i} - \pdv{}{p_i}\pdv{H}{q_i}\right) = 0,
\end{equation}
fordi de to blandede andenafledte af $H$ er identiske og trækkes fra
hinanden. Et divergensfrit strømningsfelt bevarer volumen, ganske som
en inkompressibel væskestrømning — dette er indholdet af
\eqref{eq:liouville}.

\begin{bemaerkning}
Liouvilles sætning er grundstenen for al klassisk statistisk mekanik:
den garanterer, at en sandsynlighedsfordeling over faserummet
udvikler sig som et inkompressibelt "faserumsfluidum", hvilket er
udgangspunktet for Liouville-ligningen for fordelingsfunktionen
$\rho(q,p,t)$. Vi forfølger ikke denne retning i dette kompendium, men
den er værd at nævne, fordi den samme volumenbevarelse (nu udtrykt som
unitaritet af tidsudviklingen) er en af de mest grundlæggende
egenskaber, kvantemekanikken må arve fra den klassiske teori.
\end{bemaerkning}

\begin{figure}[htbp]
  \centering
  \begin{tikzpicture}[scale=1.0]
    \draw[-{Latex}] (-0.3,0) -- (7.5,0) node[right] {$q$};
    \draw[-{Latex}] (0,-0.3) -- (0,3.3) node[above] {$p$};
    % Oprindeligt område Omega_0 (cirkulær-ish blob)
    \draw[thick,blue!60!black,fill=blue!12] plot[smooth cycle] coordinates
      {(1.0,1.6) (1.5,2.0) (2.0,1.7) (1.9,1.1) (1.4,0.9)};
    \node[blue!60!black] at (1.5,1.5) {$\Omega_0$};
    % Deformeret område Omega_t (aflangt, samme areal, forskudt)
    \draw[thick,black,fill=black!8] plot[smooth cycle] coordinates
      {(4.6,2.4) (5.6,2.55) (6.5,2.1) (6.2,1.55) (5.2,1.5) (4.5,1.85)};
    \node at (5.5,2.0) {$\Omega_t$};
    % Pil der viser strømningen/tidsudviklingen
    \draw[-{Latex},thick,dashed,green!55!black] (2.1,1.4) to[bend left=15] (4.4,1.85);
    \node[green!55!black] at (3.2,2.15) {tidsudvikling};
  \end{tikzpicture}
  \caption{Liouvilles sætning: et område $\Omega_0$ i faserummet
  deformeres under tidsudviklingen til $\Omega_t$ — formen ændres,
  men arealet (i to dimensioner) hhv.\ volumenet (i $2n$ dimensioner)
  er bevaret.}
  \label{fig:liouville}
\end{figure}

\section{Opsummering}

\begin{itemize}
  \item Poisson-parentesen $\{f,g\}=\sum_i(\partial_{q_i}f\,\partial_{p_i}g - \partial_{p_i}f\,\partial_{q_i}g)$
        opsummerer tidsudviklingen af enhver faserumsfunktion:
        $df/dt = \{f,H\} + \partial f/\partial t$.
  \item De fundamentale parenteser $\{q_i,p_j\}=\delta_{ij}$,
        $\{q_i,q_j\}=\{p_i,p_j\}=0$ definerer, hvad det betyder at være
        kanonisk konjugerede variable.
  \item Poisson-parentesen er antisymmetrisk, bilineær, opfylder
        Leibniz-reglen og Jacobi-identiteten — den er en Lie-parentes.
  \item En tidsuafhængig størrelse $f$ er bevaret netop når
        $\{f,H\}=0$; bevarede størrelser er lukkede under
        Poisson-parentesen.
  \item Banemomentets komponenter opfylder $\{L_x,L_y\}=L_z$ (cyklisk)
        — den klassiske forløber for rotationsgruppens
        kommutationsrelationer.
  \item Liouvilles sætning: faserumsvolumen er bevaret under
        Hamiltons ligningers tidsudvikling, fordi strømningsfeltet er
        divergensfrit.
\end{itemize}

\begin{opgave}
Vis direkte fra definitionen \eqref{eq:def-poisson}, at Poisson-
parentesen er antisymmetrisk og opfylder Leibniz-reglen
$\{fg,h\}=f\{g,h\}+\{f,h\}g$.
\end{opgave}

\begin{opgave}
Brug Jacobi-identiteten til at vise, at hvis $\{f,H\}=0$ og
$\{g,H\}=0$ (begge tidsuafhængige og bevarede), så er
$\{\{f,g\},H\}=0$ — altså at $\{f,g\}$ også er bevaret. (Antag
$\partial f/\partial t=\partial g/\partial t=0$.)
\end{opgave}

\begin{opgave}
Gennemfør den fulde udregning af $\{L_x,L_y\}$ fra definitionen
\eqref{eq:def-poisson}, led for led, og bekræft resultatet
\eqref{eq:Lx-Ly}. Vis derudover, at $\{L_z,x\}=-y$ og $\{L_z,y\}=x$
(dvs.\ at $L_z$ genererer infinitesimale rotationer om $z$-aksen —
en forsmag på den generelle forbindelse mellem Poisson-parenteser og
symmetritransformationer, som vender tilbage i
kapitel~\ref{kap:kanoniske-transf}).
\end{opgave}
